\documentclass[11pt]{article}
\usepackage[a4paper,margin=25mm]{geometry}
\usepackage{amsmath,amssymb,amsthm,booktabs}
\usepackage[hidelinks]{hyperref}
\emergencystretch=2em
\newtheorem{theorem}{Theorem}
\newcommand{\Z}{\mathbb Z}
\title{A nonnumerical multiplier in an odd-order difference set\\
\large Disproof of Conjecture 00000008412}
\author{gaochengzhecpu}\date{}
\begin{document}\maketitle
\begin{abstract}
We exhibit a $(27,13,6)$ difference set in $(\Z/3\Z)^3$ preserved by
the coordinate cycle $(a,b,c)\mapsto(b,c,a)$. This is a group
automorphism and a nonnumerical multiplier. The ambient group has
no even-order subgroup or subgroup quotient, and no index-two
subgroup. Thus nonnumerical multipliers are not confined to the
asserted $(16,6,2)$ exception or its index-two extension.
All finite difference counts and the automorphism are verified in Lean.
\end{abstract}

\section*{Definitions and scope}
For an additive finite group $G$, a $(v,k,\lambda)$ difference set
is a $k$-element subset $D$ with $|G|=v$ such that each nonzero
$g\in G$ occurs exactly $\lambda$ times as $x-y$ with $(x,y)\in D^2$.
A multiplier is an automorphism $\sigma$ of $G$ such that
$\sigma(D)=D+t$ for some $t\in G$. In an abelian group a multiplier
is numerical if $\sigma(x)=nx$ for an integer $n$. These are the
usual abelian conventions; see \'{O} Cath\'{a}in, Section 3.1.

We refute the only-if direction of the conjecture's second clause:
the restriction of nonnumerical multipliers to the stated even-order
exception. The source does not separately define ``coexists.''
The exclusion proved here covers an exception in the ambient group,
a subgroup, a quotient of a subgroup, or an index-two extension
of a group of order $16$. Indeed, no even-order group section is
possible. We do not assign a meaning involving unrelated groups to
that word. No claim about the two other clauses is needed.

\section*{The difference set}
Let $G=(\Z/3\Z)^3$, with coordinatewise addition and subtraction,
and let $D$ be the following set of thirteen triples:
\[
\begin{split}
D=\{&(0,0,1),(0,1,0),(0,1,1),(0,1,2),\\
    &(1,0,0),(1,0,1),(1,1,0),(1,1,1),\\
    &(1,2,0),(1,2,2),(2,0,1),(2,1,2),(2,2,1)\}.
\end{split}
\]
Its elements are distinct. For $g\in G$ write
$N(g)=\#\{(x,y)\in D^2:x-y=g\}$. Direct subtraction modulo $3$
gives the full table below. The row is the first coordinate of $g$;
the column is its last two coordinates.
\[
\begin{array}{c|rrrrrrrrr}
 &00&01&02&10&11&12&20&21&22\\\hline
0&13&6&6&6&6&6&6&6&6\\
1&6&6&6&6&6&6&6&6&6\\
2&6&6&6&6&6&6&6&6&6
\end{array}
\]
For example, for $g=(0,0,1)$ the six possible first entries $x$
are $(0,1,0)$, $(0,1,1)$, $(0,1,2)$, $(1,0,1)$, $(1,1,1)$, and $(1,2,0)$;
each has $y=x-g\in D$. Each table entry can equivalently be obtained
by testing $x-g\in D$ for the thirteen displayed $x$.
Thus $D$ is a nontrivial $(27,13,6)$ difference set.

\section*{A multiplier that is not numerical}
Define $T(a,b,c)=(b,c,a)$. This map is additive, has inverse
$T^{-1}(a,b,c)=(c,a,b)$, and satisfies $T^3=1$. The displayed set
$D$ is a union of the following $T$-orbits:
\[
\begin{gathered}
\{(1,1,1)\},\\
\{(0,0,1),(0,1,0),(1,0,0)\},\\
\{(0,1,1),(1,1,0),(1,0,1)\},\\
\{(0,1,2),(1,2,0),(2,0,1)\},\\
\{(1,2,2),(2,2,1),(2,1,2)\}.
\end{gathered}
\]
Consequently $T(D)=D$, so $T$ is a multiplier with translation $0$.
It is not numerical: for every integer $n$,
\[
 T(1,0,0)=(0,0,1)\ne(n\bmod3,0,0)=n(1,0,0).
\]
This excludes every integer scalar map, and therefore also all
numerical automorphisms with the usual coprimality condition.

\section*{Excluding the proposed exception}
The ambient group has order $27$, not $16$ or $32$. More strongly,
if $S\le G$ and $K\le S$, Lagrange's theorem gives
\[
 |S/K|\mid |S|\mid27.
\]
Hence every such group section has odd order. In particular neither
an order-$16$ group supporting a $(16,6,2)$ difference set nor its
order-$32$ index-two extension can occur in any such section.
Also $[G:S]\mid27$, so $G$ has no subgroup of index two.

\begin{theorem}
A nontrivial difference set can have a nonnumerical multiplier in
an odd-order group having no even-order group section. Thus the
claimed restriction to the $(16,6,2)$ exception or its index-two
extension in Conjecture 00000008412 is false.
\end{theorem}
\begin{proof}
The difference set $D\subseteq(\Z/3\Z)^3$ and its multiplier $T$
have exactly these properties.
\end{proof}

\section*{Formal correspondence and reproduction}
Lean uses the actual product additive group
\texttt{ZMod 3}\,$\times$\,\texttt{ZMod 3}\,$\times$\,\texttt{ZMod 3}.
The predicate \texttt{IsDifferenceSet} counts ordered pairs with
their actual group difference. The full $(27,13,6)$ property is
kernel checked. The map \texttt{cycle} is a genuine additive
equivalence with an explicit inverse; its additivity, its cube,
and preservation of $D$ are proved. The multiplier predicate allows
translations. The theorem \texttt{cycle\_not\_numerical} quantifies
over all integers, not just selected residues.

The theorem \texttt{no\_even\_section} uses Mathlib's Lagrange theorem
for arbitrary subgroups and subgroup quotients. Together with the
index-two exclusion it proves the structural facts in the final
\texttt{counterexample}. These are group properties, not tests on
a list of guessed subgroup orders.

Lean 4.19.0 and Mathlib revision\\
\texttt{c44e0c8ee63ca166450922a373c7409c5d26b00b}
are pinned. In \texttt{lean/}, run
\begin{verbatim}
lake exe cache get
lake build
lake env lean -DwarningAsError=true Main.lean
\end{verbatim}
The submitted project is rebuilt in a fresh directory. There are
no proof placeholders, custom axioms, or external computational
oracles. The accompanying \texttt{verify.py} reproduces the original
orbit-union search and all finite counts using Python's standard
library; Lean does not trust its output. Run \texttt{python verify.py}
for that check and \texttt{tectonic main.tex} to regenerate the PDF.
The author performed a separate adversarial self-review; no
independent reviewer or subagent was used.

\section*{References}
\begin{enumerate}
\item The Last Math Competition, Conjecture 00000008412. Its exact
bilingual source and upstream provenance are included in the submission.
\item P. \'{O} Cath\'{a}in, \emph{Inequivalence of difference sets:
On a remark of Baumert}, Electronic Journal of Combinatorics
20(1) (2013), P38, Section 3.1.
\href{https://www.combinatorics.org/ojs/index.php/eljc/article/download/v20i1p38/pdf/}{Publisher's full text}.
\end{enumerate}
\end{document}
